\documentclass[../main.tex]{subfiles}
\graphicspath{{\subfix{../assets/}}}

\begin{document}
\subsection{Types of Ferrous Metals}
\subsubsection{409 Stainless Steel}
\begin{table}[h!] % 140
	\begin{center}
		\begin{tabularx}{\textwidth}{
				| >{\RaggedRight\arraybackslash}X
				| >{\Centering\arraybackslash}X
				| >{\RaggedRight\arraybackslash}X |}
			\hline
			\multicolumn{3}{|c|}{\includegraphics[width=0.55\textwidth]{409ssfig1.jpg}}                                                                                                                                                                                                                                                                                                             \\
			\multicolumn{3}{|c|}{Figure 1: \textbf{409 Stainless Steel (Exhaust Pipes)}}                                                                                                                                                                                                                                                                                                            \\
			\multicolumn{3}{|c|}{\hspace*{\fill}{\scriptsize \parencite{ss409pic}\quad}}                                                                                                                                                                                                                                                                                                            \\
			\hline
			Applications                                                                                                                                                                                                                                                                                                       & Nominal Composition (S40910) & Properties \& Justification for use \\
			\hline
			409 stainless steel is commonly used for automotive exhausts but is still considered a general-purpose stainless steel. Though inferior to 304 stainless, it's suitable for applications requiring heat resistance and durability, sacrificing aesthetics for performance. It's also weldable. \parencite{ss409ds} &
			$\sim0.08\%\mathrm{C}$, $\sim1\%\mathrm{Mn}$, $\sim1\%\mathrm{Si}$, $\sim0.045\%\mathrm{Si}$, $\sim0.03\%\mathrm{S}$, $\sim11.1\%\mathrm{Cr}$, $\sim0.5\%\mathrm{Ni}$, $\sim0.6\%\mathrm{Ti}$, $\sim85.645\%\mathrm{Fe}$  \parencite{ss409ds}                                                                      &
			The extreme temperatures exhaust systems are exposed to during regular use necessitates a corrosion and heat-resistant steel. While aesthetically less desirable due to surface oxidation, the increased service life relative to its cost in comparison to more expensive examples \textit{i.e.} 304 stainless, more than justify its use. \parencite{ss409vs304}                      \\
			\hline
		\end{tabularx}
	\end{center}
\end{table}
\newpage

\subsubsection{Cast Iron}
\begin{table}[h!] % 157
	\begin{center}
		\begin{tabularx}{\textwidth}{
				| >{\RaggedRight\arraybackslash}X
				| >{\Centering\arraybackslash}X
				| >{\RaggedRight\arraybackslash}X |}
			\hline
			\multicolumn{3}{|c|}{\includegraphics[width=0.55\textwidth]{castironfig1.jpg}}                                                                                                                                                                                                                                                                                                                                            \\
			\multicolumn{3}{|c|}{Figure 2: \textbf{Cast Iron (Engine Block)}}                                                                                                                                                                                                                                                                                                                                                         \\
			\multicolumn{3}{|c|}{\hspace*{\fill}{\scriptsize \parencite{castironpic}\quad}}                                                                                                                                                                                                                                                                                                                                           \\
			\hline
			Applications                                                                                                                                                                                                                                                                                                                                & Nominal Composition (White Cast Iron) & Properties \& Justification for use \\
			\hline
			The most obvious use of cast iron parts is in engine blocks and gearbox cases, where its ease of manufacture as well as excellent durability lends itself to see usage even into the age of engineered composites. Cast iron isn't limited to gearboxes, however, seeing extensive use in braking systems as well. \parencite{castironapps} &
			$\sim2.4\%\mathrm{C}$, $\sim0.525\%\mathrm{Mn}$, $\sim0.2\%\mathrm{S}$, $\sim1.3\%\mathrm{Si}$, $\sim0.125\%\mathrm{P}$, $\sim95.45\%\mathrm{Fe}$ \parencite{castironcomp}                                                                                                                                                                  &
			Cast iron's popularity is due to its mechanical properties, low cost, and ease of manufacture. \textcite{engtextbookcastiron} highlights its excellent properties, including low cost, high compressive strength and rigidity, easy machinability, and high fluidity. These properties make cast iron a convenient material for its affordability.                                                                        \\
			\hline
		\end{tabularx}
	\end{center}
\end{table}
\newpage

\subsubsection{Low-carbon Steel}
\begin{table}[h!] % 200
	\begin{center}
		\begin{tabularx}{\textwidth}{
				| >{\RaggedRight\arraybackslash}X
				| >{\Centering\arraybackslash}X
				| >{\RaggedRight\arraybackslash}X |}
			\hline
			\multicolumn{3}{|c|}{\includegraphics[width=0.55\textwidth]{car_steel_panels.png}}                                                                                                                                                                                                                                                                                                                                                                                                                                                      \\
			\multicolumn{3}{|c|}{Figure 3: \textbf{Low-carbon (Mild) Steel}}                                                                                                                                                                                                                                                                                                                                                                                                                                                                        \\
			\multicolumn{3}{|c|}{\hspace*{\fill}{\scriptsize \parencite{mildsteelpic}\quad}}                                                                                                                                                                                                                                                                                                                                                                                                                                                        \\
			\hline
			Applications                                                                                                                                                                                                                                                                                                                                                                                                                                      & Nominal Composition (AISI 1010) & Properties \& Justification for use               \\
			\hline
			Mild steel is often the material of choice for nonstructural automotive components, such as crumple zones and impact-deforming panels. Low-carbon steel is unique in that it deforms upon impact, unlike reinforced steels that either don't deform or shatter. Mild steel is also easily recyclable, making it suitable for sustainable design. Its affordability and widespread availability make it an ideal solution for nonstructural parts. &
			$\sim0.105\%\mathrm{C}$, \hskip 3cm $\sim0.4\%\mathrm{Mn}$, $\sim0.04\%\mathrm{P}$, $\sim0.05\%\mathrm{S}$, $\sim99.405\%\mathrm{Fe}$ \parencite{mildsteelcompprop}                                                                                                                                                                                                                                                                               &
			It is precisely the mechanical properties that make low-carbon steel undesirable for other applications, that crown it the \textit{de facto} material of choice for body components. It's necessary for the crash structure to deform upon impact to reduce the energy transferred to any occupants; hardened tool steels or cast iron, both brittle and highly dense, are unsuitable for this task. The mechanical ability of low-carbon steel to deform freely on impact is invaluable in an automotive. \parencite{mildsteelappprop} \\
			\hline
		\end{tabularx}
	\end{center}
\end{table}
\newpage

\subsubsection{High-strength-low-alloy Steel}
\begin{table}[h!] % 164
	\begin{center}
		\begin{tabularx}{\textwidth}{
				| >{\RaggedRight\arraybackslash}X
				| >{\Centering\arraybackslash}X
				| >{\RaggedRight\arraybackslash}X |}
			\hline
			\multicolumn{3}{|c|}{\includegraphics[width=0.55\textwidth]{hslapic.jpg}}                                                                                                                                                                                                                                                                                                                                                                                                                                                        \\
			\multicolumn{3}{|c|}{Figure 4: \textbf{High-strength-low-alloy (HSLA) Steel}}                                                                                                                                                                                                                                                                                                                                                                                                                                                    \\
			\multicolumn{3}{|c|}{\hspace*{\fill}{\scriptsize \parencite{hslapic}\quad}}                                                                                                                                                                                                                                                                                                                                                                                                                                                      \\
			\hline
			Applications                                                                                                                                                                                                                                                                                                                                                                                                                                                         & Nominal Composition & Properties \& Justification for use \\
			\hline
			HSLA steels regularly find use in automotive chassis, their superior strength-to-weight ratio extracting extra performance while remaining safe. \parencite{hslaprosconsapp} Specifically, bumpers, suspension systems, and structural panels are typically constructed of HSLA steel, noting that these parts are both stronger and lighter than alternatives.                                                                                                      &
			Unfortunately, due to the greatly varied nature of HSLA steels, the following nominal composition is only a reference and is likely not representative of any individual example: $\sim0.25\%\mathrm{C}$, $\sim2\%\mathrm{Mn}$, \hskip 3cm $\sim97\%\mathrm{Fe}$. The remaining weight is made up of some combination of: $\mathrm{V}$, $\mathrm{Nb}$, $\mathrm{Cr}$, $\mathrm{N}$, $\mathrm{Mo}$, $\mathrm{Ni}$, $\mathrm{Cu}$, $\mathrm{Zr}$. \parencite{hslacomp} &
			The "exceptional strength and toughness" of HSLA steels under high-intensity applications is ideal for materials in the automotive industry. \parencite{hslaprosconsapp} Not only is HSLA steel remarkably durable, its (relative) freedom in alloying composition enables material engineers to specify properties down to a microscopic scale, incentivising targeted solutions via its excellent customisability.                                                                                                             \\
			\hline
		\end{tabularx}
	\end{center}
\end{table}
\newpage

\subsection{Surface Treatments of Ferrous Metals}
\subsubsection{Nitriding of Engine Components} % 200
Steels found in engine components are exposed to extreme mechanical loads and temperatures; the crankshaft is a common example of nitrided steel, necessary for its additional surface hardness to mitigate wear while retaining the mechanical properties of the substrate.
As a case-hardening process, nitrdiding infuses nitrogen into the surface of the steel, forming a superficial compound layer with superior hardness. \parencite{nitriding2}
This process is best explained with a diagram.

\setcounter{figure}{4}
\begin{figure}[h!]
	\begin{center}
		\includegraphics[width=0.75\textwidth, angle=-0.8]{gas-nitriding.jpg}
	\end{center}
	\caption{The gas nitriding process \parencite{nitriding1}}
\end{figure}

At a basic level, the steps of (gas) nitriding are as follows:
\begin{enumerate}
	\item The part to be nitrided is placed in a furnace
	\item The atmosphere of the furnace is replaced with anhydrous ammonia gas (${\mathrm{NH}_3}_{(\text{g})}$)
	\item The temperature is increased until the ammonia decomposes into nitrogen gas (${\mathrm{N}_2}_{(\text{g})}$) and diffuses into a nitride layer chemically integrated into the steel
	\item The part is held in this state until the desired nitride layer thickness is reached
\end{enumerate}

The specific outcomes of nitriding include excellent surface hardness while maintaining core ductility \parencite{nitriding2}, as well as increased chemical (corrosion) resistance, both desirable for engine components due to the frequent shock loads that they may experience.
Not only does the additional chemical resistance protect from oxidation stemming from thermal cycling, but the surface hardness also reduces frictional wear; this is highly beneficial on contact components \textit{e.g.} the crankshaft.
\newpage

\subsubsection{Electroplating of Trim Panels} % 210
While the mechanical properties of a part may be the biggest consideration to an engineer, the average customer is likely much more concerned with the aesthetic appeal of the vehicle.
Electroplating is one method suitable for achieving these aesthetic outcomes, where any desired appearances can be easily bonded as a superficial coating which additionally protects the substrate.
A standard setup for electroplating is shown in the following diagram.

\begin{figure}[h!]
	\begin{center}
		\includegraphics[width=0.75\textwidth]{electroplating_pic.jpg}
	\end{center}
	\caption{A standard electroplating cell \parencite{electroplatepic}}
\end{figure}

Electroplating utilises controlled electrolysis to coat the substrate with a pure layer of metal. Common options include copper, chromium, zinc, and nickel, but any metal available as an ion (solution) can be utilised. \parencite{electroplating}
The above diagram illustrates a scalable setup, where the dissolved ions are driven to the cathode (substrate) via an applied current; the process is largely identical regardless of part size.
In automotives, electroplating is often employed to give parts a metallic finish not only for aesthetics, but also protection to maintain the longevity of the components, allowing for aesthetic specifications which may be otherwise impossible or highly impractical \textit{e.g.} utilising pure gold trim for bumpers, etc.
Aesthetically pleasing visuals are a crucial element of automotive design, and electroplated parts are one such way to fulfil this need.
\newpage

\subsubsection{(Spray) Painting of Body Panels} % 100
Practically speaking, almost all cars in existence are (spray) painted in some capacity.
Painted surfaces offer both aesthetic appeal and corrosion resistance, both valuable properties in automotive contexts.
These paints may also be engineered to exact specifications, allowing for additional design flexibility.
This process is usually automated, as shown below.

\begin{figure}[h!]
	\begin{center}
		\includegraphics[width=0.75\textwidth]{robotspray.jpg}
	\end{center}
	\caption{A robotic spray booth \parencite{spraypic}}
\end{figure}

Some benefits of painting parts include:
\begin{itemize}
	\item Resistance to corrosion
	\item Aesthetic appeal
	\item Protection against environmental wear
\end{itemize}
It's necessary for parts exposed to the elements in routine use to be shielded from degradation and damage, and painting provides an optimal solution that fulfils both these requirements.
\newpage

\subsubsection{Carburising of Gearbox Components} % 160
Similar to nitriding, carburisation is another case-hardening process utilised in similar applications.
In automotives, gearbox components are a common example of carburised steel, where durable contact surfaces are necessary, but brittle materials are unacceptable.
Carburisation provides an alternative to nitriding in this case, though it is less practically effective and efficient, requiring higher temperatures but unable to achieve the same level of surface hardness.

\begin{figure}[h!]
	\begin{center}
		\includegraphics[width=0.75\textwidth]{carburising.jpg}
	\end{center}
	\caption{A carburising furnace \parencite{carbpic}}
\end{figure}

The process of carburising is similar to nitriding:
\begin{enumerate}
	\item The part is held in a furnace with hydrocarbon gases and heated
	\item The hydrocarbons decompose and the free carbon diffuses into the steel
	\item The part is quenched to form a martensitic layer and lock in the carbon
\end{enumerate}

\textcite{carbdesc} outlines the widespread application of carburising in manufacturing gears; contact surfaces within gearboxes require extremely high surface hardness, but must maintain a degree of elasticity to provide a margin of safety in cases of shock loading.
Carburisation fulfils this similarly to nitriding, leaving a hard case but ductile interior.
\newpage

\sloppy
\printbibliography[heading=subbibnumbered,keyword={sect1},title={References}]
\fussy
\end{document}
